\par\Needspace{6\baselineskip}\medskip
\begin{center}\hypertarget{native-174-article127}{}\pdfbookmark[2]{127.}{article-127}127.\end{center}
Postquam tres loci heliocentrici eo quem in art. praec. descripsimus modo
eruti sunt, sequenti quoque modo procedi poterit. Determinentur elementa
reliqua per problema in artt. 85—105 tractatum primo e loco primo et secundo
cum intervallo temporis respondente, dein vero eodem modo e loco secundo
et tertio temporisque intervallo respondente: ita pro singulis elementis
duo valores prodibunt, e quorum differentiis duas ad libitum pro $X$ et $Y$
accipere licebit. Magnopere hanc methodum commendat commodum haud
spernendum, quod in hypothesibus primis elementa reliqua, praeter duo ea
quae ad stabiliendum $X$ et $Y$ eliguntur, omnino negligere licet, quae in ultimo
demum calculo, valoribus correctis ipsarum $x,y$ superstructo, determinabuntur
sive e sola combinatione prima, sive e sola secunda, sive quod plerumque
praeferendum est e combinatione loci primi cum tertio. Ceterum electio illorum
duorum elementorum, quae generaliter loquendo arbitraria est, magnam
solutionum varietatem suppeditat; adoptari poterunt e. g. logarithmus semiparametri
cum logarithmo semiaxis maioris, vel prior cum excentricitate, vel
cum eadem posterior, vel cum aliquo horum elementorum longitudo perihelii:
combinari quoque poterit aliquod horum quatuor elementorum cum anomalia
excentrica loco medio in utroque calculo respondente, siquidem orbita elliptica
evaserit, ubi formulae 27—30 art. 96 calculum maxime expeditum afferent.
In casibus specialibus autem haec electio quadam circumspectione indiget;
ita e. g. in orbitis ad parabolae similitudinem vergentibus semiaxis maior $a$
ipsiusve logarithmus minus idonei forent, quippe quorum variationes immodicae
variationibus ipsarum $x,y$ haud proportionales censeri possent: in tali
casu magis e re esset eligere $\dfrac{1}{a}\mathpunct{\cdot}$ Sed his cautelis eo minus immoramur,
quum methodus quinta in art. seq. explicanda quatuor hactenus expositis in
omnibus fere casibus palmam praeripiat.
